% Generated by scripts/build_cv_variants.py from cv-latex/variants.json.
% Edit the bilingual content source, then regenerate.
\documentclass[10pt,a4paper]{article}

\usepackage[a4paper,top=10mm,bottom=9mm,left=13mm,right=13mm,footskip=8mm]{geometry}
\usepackage{fontspec}
\usepackage[AutoFakeBold=false]{xeCJK}
\usepackage{xcolor}
\usepackage[protrusion=true]{microtype}
\usepackage{ragged2e}
\usepackage{enumitem}
\usepackage{titlesec}
\usepackage{tabularx}
\usepackage{array}
\usepackage{needspace}
\usepackage{fancyhdr}
\usepackage{lastpage}
\usepackage{hyperref}

\definecolor{ink}{HTML}{111111}
\definecolor{inkii}{HTML}{3A3A3A}
\definecolor{inkiii}{HTML}{6A6A6A}
\definecolor{rule}{HTML}{8C8C8C}

% Fonts are found through OPENTYPEFONTS, which scripts/compile_cv.py points
% at cv-latex/fonts/.
\setmainfont{SourceSerif4SmText-Regular.otf}[
  BoldFont=SourceSerif4SmText-Bold.otf,
  Numbers=Lining,
  ItalicFont=SourceSerif4SmText-Regular.otf]
\newfontfamily\namefont{SourceSerif4Subhead-Bold.otf}[Numbers=Lining,BoldFont=SourceSerif4Subhead-Bold.otf]
\setCJKmainfont{Noto Serif CJK SC Medium}[BoldFont={Noto Serif CJK SC Bold}]
\xeCJKsetup{CJKecglue={\hskip 0.2em plus 0.08em minus 0.04em}}
\newcommand{\spaced}[1]{\xeCJKsetup{CJKglue={\hskip 0.14em}}#1}

\hypersetup{
  colorlinks=true,
  linkcolor=ink,
  urlcolor=ink,
  pdftitle={胡力中 - 升学简历},
  pdfauthor={胡力中},
  pdfsubject={升学简历},
  pdfcreator={XeLaTeX}
}

\renewcommand{\normalsize}{\fontsize{9.8}{13.31}\selectfont}
\renewcommand{\small}{\fontsize{8.62}{11.71}\selectfont}
\renewcommand{\footnotesize}{\fontsize{8.62}{11.71}\selectfont}
\normalsize
\AtBeginDocument{\color{ink}}

\setlength{\parindent}{0pt}
\setlength{\parskip}{0pt}
\setlist[itemize]{leftmargin=1.05em,labelsep=0.5em,itemsep=0.4pt,topsep=1pt,parsep=0pt,partopsep=0pt,
  label={\raisebox{0.2em}{\rule{0.28em}{0.28em}}}}

% Section headings: bold, spaced capitals (or spaced Chinese), a grey hairline.
\titleformat{\section}{\fontsize{9.8}{13.31}\selectfont\bfseries}{}{0pt}{\spaced}
  [{\vspace{1pt}\color{rule}\titlerule[0.4pt]}]
\titlespacing*{\section}{0pt}{4pt}{2pt}
\renewcommand{\arraystretch}{1.02}
\newcolumntype{Y}{>{\raggedright\arraybackslash}X}

% Quiet secondary text in place of italics (the site uses no italics).
\newcommand{\subtle}[1]{\leavevmode{\color{inkii}#1}}

% Entry heading: title in bold, date in grey on the right.
\newlength{\entrydatewidth}
\newcommand{\entry}[2]{%
  \par\addvspace{3pt}\noindent
  \if\relax\detokenize{#2}\relax
    \textbf{#1}\par
  \else
    \settowidth{\entrydatewidth}{{\small #2}}\addtolength{\entrydatewidth}{1.5pt}%
    \begin{tabularx}{\textwidth}{@{}Y@{\hspace{1.2em}}>{\raggedleft\arraybackslash}p{\entrydatewidth}@{}}
      \textbf{#1} & \leavevmode{\small\color{inkii}#2}
    \end{tabularx}\par
  \fi
  \nobreak\vspace{-0.5pt}
}

\newcommand{\skillrow}[2]{%
  \par\noindent\hangindent=1em\hangafter=1\textbf{#1：}#2\par
}

\pagestyle{empty}

\begin{document}
\RaggedRight
\setlength{\emergencystretch}{1em}
{\namefont\bfseries\fontsize{20}{22}\selectfont 胡力中}\hspace{0.6em}{\fontsize{13}{16}\selectfont\color{inkiii}Lizhong Hu}\par
\vspace{2pt}{\small\color{inkii}\href{https://au1bhi.com}{au1bhi.com}\enspace\textperiodcentered\enspace \href{mailto:admin@au1bhi.com}{admin@au1bhi.com}\enspace\textperiodcentered\enspace \href{https://github.com/au1bhi}{GitHub}\enspace\textperiodcentered\enspace \href{https://codeforces.com/profile/Au1Bhi}{Codeforces}\enspace\textperiodcentered\enspace \href{https://www.linkedin.com/in/lizhong-hu-686b9b34a}{LinkedIn}}\par
\vspace{3pt}{\small\bfseries 信息检索 · 可验证 RAG · 检索评测}\par
\vspace{3pt}{\color{ink}\rule{\textwidth}{0.7pt}}\par
\vspace{-8pt}
\section{个人简介}
中美 4+0 双学位软件工程专业本科生，预计 2027 年毕业。毕业设计为可验证引用的 RAG 问答系统 NoteLLM，研究兴趣为信息检索、可验证 RAG 与可靠软件系统；有后端开发和程序设计竞赛经验。计划攻读计算机相关专业研究生。

\section{教育经历}
\begin{tabularx}{\textwidth}{@{}Y>{\raggedleft\arraybackslash}p{0.24\textwidth}@{}}
  \textbf{泉州信息工程学院} & \subtle{2023--2027（预计）} \\
  软件工程（中外合作办学） & GPA: 3.10/4.00 \\[0.5pt]
  \multicolumn{2}{@{}l@{}}{\textbf{美国宾州滑石大学（Slippery Rock University）}} \\
  信息系统（Information Systems），理学学士（在读，预计 2027） & GPA: 3.583/4.000
\end{tabularx}
\par\vspace{2pt}{\footnotesize\color{inkiii}中美 4+0 双学位联合培养本科项目，全程在中国完成培养，经中华人民共和国教育部正式批准认可。}
\section{研究方向与兴趣}
算法与数据结构；软件与分布式系统；信息检索与检索增强生成；可靠智能信息系统。
\section{毕业设计与项目}
\entry{\href{https://github.com/au1bhi/NoteLLM}{NoteLLM：基于可验证 RAG 的个人学习问答系统}}{进行中}
\subtle{本科毕业设计}
\begin{itemize}
  \item 基于 FastAPI、异步 SQLAlchemy、PostgreSQL 与 pgvector 构建后端，实现余弦相似度检索；支持 PDF、Markdown、TXT 的按页提取与重叠分块。
  \item 在服务端校验引用是否来自检索分块并保存原文摘录，支持问答结果的来源核查。
  \item 在 7 份合成资料、34 道固定问题的单次基线测试中，来源级 Recall@5 为 100\%，自动引用来源匹配率为 97.1\%，平均检索耗时 339 ms。以上衡量该合成测试集的检索与引用来源匹配，不代表答案准确率或通用场景性能。
  \item 评测报告与源码公开：github.com/au1bhi/NoteLLM；评测结果位于 docs/evaluation/latest-results.md。
\end{itemize}
\entry{e体+（ETIPlus）：社区治理与微信群管理平台后端}{2025 年 3–4 月}
\begin{itemize}
  \item 使用 Flask、PostgreSQL 与 MongoDB 实现信用赋分、异常审计及群管理 API；通过 Postman 执行回归检查。
\end{itemize}
\section{命题与竞赛组织}
\entry{\href{https://www.matiji.net/exam/noi}{码蹄集 NOI-Pre 题库｜志愿验题员}}{2024 年 7–10 月}
\begin{itemize}
  \item 参与码蹄集与 CCF 合作题库测试，构造边界用例，通过定向提交验证题解正确性，排查并反馈测试数据与判题结果缺陷。 \allowbreak\enspace{\small \mbox{\href{https://www.matiji.net/exam/volunteer}{志愿者成员展示}}}
\end{itemize}
\entry{\href{https://new.saikr.com/vse/adpc/2025/spring}{第六届大学生算法设计与编程挑战赛（春季赛）}}{2025 年春季}
\subtle{命题组成员}
\begin{itemize}
  \item 参与题目设计、题解验证、难度校准与题解审阅。
\end{itemize}
\entry{校 ICPC/CCPC 集训队｜协调人}{2024 年 9 月–2026 年 9 月}
\begin{itemize}
  \item 协调校 ICPC/CCPC 集训队招新、训练与队伍选拔，组织本校队伍参赛事务。
  \item 组织本校队伍参加：2026 ICPC 全国邀请赛（南昌）、2025 CCPC 全国邀请赛（南昌）暨江西省赛、2026 CCPC 全国邀请赛（福建）暨福建省赛、2025 ICPC Asia EC 网络预选赛与 CCPC 网络赛、2026 CCPC 贵州省赛暨全国邀请赛、2025 ICPC 贵州省赛、2024 ICPC 昆明区域赛。
\end{itemize}
\entry{校内程序设计竞赛组织与命题}{}
\begin{itemize}
  \item “国教AC杯”大学生程序设计竞赛：主要组织者兼命题人。
  \item 2025 泉州信息工程学院软件学院第二届“爱拼会赢”科创节 ACM 算法比赛：参与组织与命题。
\end{itemize}
\section{程序设计竞赛与获奖}
\begin{itemize}
  \item \href{https://au1bhi.com/files/WorldQuant-Challenge-Silver-Certificate.pdf}{WorldQuant BRAIN 量化挑战赛（WorldQuant Challenge）：Silver Level（银级），2026 年 11 月。}
  \item 2026 中国大学生程序设计竞赛全国邀请赛（福建）暨第十三届 CCPC 福建省大学生程序设计竞赛：最快解题奖（A 题，单题奖项）。
  \item 第49届 ICPC 国际大学生程序设计竞赛亚洲区域赛（昆明）：Honorable Medal，2024。
  \item 第十五届蓝桥杯全国软件和信息技术专业人才大赛：福建省赛 Python 程序设计大学 B 组一等奖，2024。
\end{itemize}
\section{技术技能}
\skillrow{编程与算法}{C++、Python、Java、JavaScript、Go；算法与数据结构。}
\skillrow{后端与数据库}{FastAPI、Flask、RuoYi、PostgreSQL、MongoDB、MySQL；REST API 设计与系统集成。}
\skillrow{检索与 AI 系统}{RAG、大语言模型 API 集成、Prompt 设计。}
\skillrow{系统与运维}{Linux、Docker、Git、DOMjudge、Cloudflare Workers/CDN/DNS、WARP 出口路由。}
\skillrow{命题与评测}{熟悉 Polygon（Codeforces）与牛客出题平台；题目准备、测试数据管理和题解验证。}
\end{document}
